% Festlegung des Allgemeinen Dokumentenformats
\documentclass[a4paper,12pt,headsepline]{scrartcl}

% Umlaute unter UTF8 nutzen
\usepackage[utf8]{inputenc}

% Variablen
%Variablen welche innerhalb der gesamten Arbeit zur Verfügung stehen sollen
\newcommand{\titleDocument}{Bachelorarbeit}
\newcommand{\subjectDocument}{
    Analyse der Echtzeitfähigkeit von Micro-ROS und FreeRTOS am Beispiel einer
    Robotersteuerungssoftware
}


% weitere Pakete
% Grafiken aus PNG Dateien einbinden
\usepackage{graphicx}
\usepackage{tikz}

% Deutsche Sonderzeichen und Silbentrennung nutzen
\usepackage[ngerman]{babel}

% Eurozeichen einbinden
\usepackage[right]{eurosym}

% Zeichenencoding
\usepackage[T1]{fontenc}

\usepackage{lmodern}

% floatende Bilder ermöglichen
%\usepackage{floatflt}

% mehrseitige Tabellen ermöglichen
\usepackage{longtable}

% Unterstützung für Schriftarten
%\newcommand{\changefont}[3]{ 
%\fontfamily{#1} \fontseries{#2} \fontshape{#3} \selectfont}

% Packet für Seitenrandabständex und Einstellung für Seitenränder
\usepackage{geometry}
\geometry{left=3.5cm, right=2cm, top=2.5cm, bottom=2cm}

% Paket für Boxen im Text
\usepackage{fancybox}

% bricht lange URLs "schön" um
\usepackage[hyphens,obeyspaces,spaces]{url}

% Paket für Textfarben
\usepackage{color}

% Mathematische Symbole importieren
\usepackage{amssymb}

\usepackage{amsmath}

% auf jeder Seite eine Überschrift (alt, zentriert)
%\pagestyle{headings}

\usepackage{pdfpages}

% erzeugt Inhaltsverzeichnis mit Querverweisen zu den Abschnitten (PDF Version)
\usepackage[bookmarksnumbered,pdftitle={\titleDocument},hyperfootnotes=false]{hyperref}
%\hypersetup{colorlinks, citecolor=red, linkcolor=blue, urlcolor=black}
%\hypersetup{colorlinks, citecolor=black, linkcolor= black, urlcolor=black}

% neue Kopfzeilen mit fancypaket
\usepackage{fancyhdr} %Paket laden
\pagestyle{fancy} %eigener Seitenstil
\fancyhf{} %alle Kopf- und Fußzeilenfelder bereinigen
\fancyhead[L]{\nouppercase{\leftmark}} %Kopfzeile links
\fancyhead[C]{} %zentrierte Kopfzeile
\fancyhead[R]{\thepage} %Kopfzeile rechts
\renewcommand{\headrulewidth}{0.4pt} %obere Trennlinie
%\fancyfoot[C]{\thepage} %Seitennummer
%\renewcommand{\footrulewidth}{0.4pt} %untere Trennlinie

% change font to serif for headings - buggy
% \usepackage{sectsty}
% \allsectionsfont{\normalfont\bfseries}

% für Tabellen
\usepackage{array}

% Runde Klammern für Zitate
%\usepackage[numbers,round]{natbib}

% Festlegung Art der Zitierung - Havardmethode: Abkuerzung Autor + Jahr
\bibliographystyle{alphadin}

% Schaltet den zusätzlichen Zwischenraum ab, den LaTeX normalerweise nach einem Satzzeichen einfügt.
%\frenchspacing

% Paket für Zeilenabstand
\usepackage{setspace}

% für Bildbezeichner
\usepackage{capt-of}

% für Stichwortverzeichnis
\usepackage{makeidx}

%Konfiguriere das Inhaltsverzeichnis
\usepackage{tocbasic}
\DeclareTOCStyleEntries[
  raggedentrytext,
  numwidth=0pt,
  numsep=1ex,
  dynnumwidth,
]{tocline}{chapter,section,subsection,subsubsection,paragraph,subparagraph}
\DeclareTOCStyleEntries[
  indent=0pt,
  linefill=\TOCLineLeaderFill,
]{tocline}{section,subsection,subsubsection,paragraph,subparagraph}


% für Listings
% \usepackage{listings}
% \lstset{numbers=left, numberstyle=\tiny, numbersep=5pt, keywordstyle=\color{black}\bfseries, stringstyle=\ttfamily,showstringspaces=false,basicstyle=\footnotesize,captionpos=b}
\usepackage{xcolor}
\usepackage[newfloat]{minted}
\usepackage{caption}
\usepackage{fancyhdr}
\newenvironment{code}{\captionsetup{type=listing}}{}
\SetupFloatingEnvironment{listing}{name=Quellcode}
\setminted{
    linenos,
	frame=single,
    bgcolor=black!5,
    fontsize=\footnotesize
}
\usepackage{graphicx}  % Add this to your preamble

\usepackage[withpage]{acronym}

% Absatz ohne Einrückung
\setlength{\parskip}{1em} % 1em entspricht der Breite eines 'M'-Zeichens
\setlength{\parindent}{0pt}

% Indexerstellung
\makeindex

% Abkürzungsverzeichnis
\usepackage[german]{nomencl}
\let\abbrev\nomenclature

% Abkürzungsverzeichnis LiveTex Version
% Titel des Abkürzungsverzeichnisses
\renewcommand{\nomname}{Abkürzungsverzeichnis}
% Abstand zwischen Abkürzung und Erläuterung
\setlength{\nomlabelwidth}{.25\textwidth}
% Zwischenraum zwischen Abkürzung und Erläuterung mit Punkten
\renewcommand{\nomlabel}[1]{#1 \dotfill}
% Variation des Abstandes der einzelnen Abkürzungen zu einander
\setlength{\nomitemsep}{-\parsep}
% Index mit Abkürzungen erzeugen
\makenomenclature
%\makeglossary

% Abkürzungsverzeichnis TeTEX Version
% \usepackage[german]{nomencl}
% \makenomenclature
% %\makeglossary
% \renewcommand{\nomname}{Abkürzungsverzeichnis}
% \AtBeginDocument{\setlength{\nomlabelwidth}{.25\columnwidth}}
% \renewcommand{\nomlabel}[1]{#1 \dotfill}
% \setlength{\nomitemsep}{-\parsep}

% Optional: Einzelne Zeilen am Anfang einer Seite unterdrücken (Schusterjungen)
% \clubpenalty = 10000
% Optional: Einzelne Zeilen am Ende einer Seite unterdrücken (Hurenkinder)
% \widowpenalty = 10000
% \displaywidowpenalty = 10000

\begin{document}
% hier werden die Trennvorschläge inkludiert
%hier müssen alle Wörter rein, welche Latex von sich auch nicht korrekt trennt bzw. bei denen man die genaue Trennung vorgeben möchte
\hyphenation{
Film-pro-du-zen-ten
Lux-em-burg
Soft-ware-bau-steins
zeit-in-ten-siv
}


% Schriftart Helvetica verwenden
%\usepackage{helvet}
%\renewcommand\familydefault{\sfdefault}

% Leere Seite am Anfang
%\thispagestyle{empty} % erzeugt Seite ohne Kopf- / Fusszeile
%\mbox{}
%\newpage

% Titelseite %
\thispagestyle{empty}

% % Logo
% \begin{figure}[t]
%  \centering
%  \includegraphics[width=0.3\textwidth]{assets/fhlogo.pdf}
% \end{figure}

\begin{center}
\textbf{\Large{~\titleDocument}}
\end{center}
\begin{verbatim}





\end{verbatim}
\begin{center}
\doublespacing
\textbf{\LARGE{\subjectDocument}}\\
\singlespacing
\end{center}
\begin{verbatim}
















\end{verbatim}
\begin{center}
\textbf{
    An der Fachhochschule Dortmund \\
    im Fachbereich Informatik \\
    Studiengang Technische Informatik \\
    erstellte Thesis \\
    zur Erlangung des akademischen Grades \\
    Bachelor of Science \\
    B. Sc. \\
}
\end{center}
\begin{center}
    \textbf{Xu, Zijian \linebreak
    geboren am 25.09.1998 \linebreak
    7204211
    }
\end{center}
\begin{flushleft}
\begin{tabular}{llll}
\textbf{Betreuung durch:} & & Prof. Dr. Christof Röhrig &\\
    & & M. Sc. Alexander Miller &\\
\textbf{Version vom:} & & Dortmund, \today &\\
\end{tabular}
\end{flushleft}


% römische Numerierung
\pagenumbering{roman}

% 1.5 facher Zeilenabstand
\onehalfspacing

\newpage

% Einleitung / Abstract
\thispagestyle{empty}
\section*{Kurzfassung}

Diese Arbeit analysiert die Echtzeitfähigkeit von Micro-ROS und FreeRTOS anhand
einer Robotersteuerungssoftware. Ziel ist der Vergleich beider Systeme
hinsichtlich ihres Echtzeitverhaltens mit Schwerpunkt auf den Ausführungszeiten.
Dazu wird zunächst die bestehende Micro-ROS-Steuerungssoftware auf FreeRTOS
portiert. Anschließend wird eine zyklengenaue Messmethode für den Programmlauf
basierend auf der Data Watchpoint and Trace Unit (DWT) der ARM-Architektur
entwickelt. Zum Abschluss werden die Laufzeitdaten visualisiert und ausgewertet.
Die Ergebnisse verdeutlichen die unterschiedlichen Schwerpunkte beider
Plattformen: Micro-ROS punktet durch die Kompatibilität mit dem ROS-Ökosystem,
während FreeRTOS mit minimalen Latenzen und deterministischem Scheduling
deutlich besseres Echtzeitverhalten bietet. Die Analyse bestätigt zudem den
maßgeblichen Einfluss der Cache-Nutzung auf die Rechenleistung.

\section*{Abstract}

This thesis analyzes the real-time capabilities of micro-ROS and FreeRTOS using
a robotic control software. The goal is to compare both systems regarding their
real-time behavior emphasized on the execution times. First, the existing
micro-ROS control software is ported to FreeRTOS. Next, a cycle-accurate
measurement method for program execution based on ARM's Data Watchpoint and
Trace Unit (DWT)is developed. Finally, the runtime data is evaluated. The
evaluation of the results highlights the different strengths of both platforms:
micro-ROS stands out due to its compatibility with the ROS ecosystem, while
FreeRTOS achieves superior real-time behavior through minimal latencies and
deterministic scheduling. Additionally, the analysis confirms the significant
impact of the cache utilization on system performance.


% einfacher Zeilenabstand
\singlespacing

\newpage
% Seitenzählung bei Inhaltsverzeichnis beginnen
\setcounter{page}{1}

% Inhaltsverzeichnis anzeigen
\thispagestyle{empty}
\tableofcontents

\newpage
% das Abbildungsverzeichnis
% Verion 1: Abbildungsverzeichnis MIT führender Nummberierung endgueltig anzeigen
\listoffigures
% Abbildungsverzeichnis soll im Inhaltsverzeichnis auftauchen
\addcontentsline{toc}{section}{Abbildungsverzeichnis}

% Verion 2: Abbildungsverzeichnis OHNE führende Nummberierung endgueltig anzeigen
%\begingroup
%\renewcommand\numberline[1]{}
%\listoffigures
%\endgroup


% % das Tabellenverzeichnis
% \newpage
% % \fancyhead[L]{Abbildungsverzeichnis / Abkürzungsverzeichnis} %Kopfzeile links
% % Tabellenverzeichnis endgültig anzeigen
% \listoftables
% % Tabellenverzeichnis soll im Inhaltsverzeichnis auftauchen
% \addcontentsline{toc}{section}{Tabellenverzeichnis}

% das Quellcodeverzeichnis
\newpage
\renewcommand*{\listlistingname}{Quellcodeverzeichnis}
\listoflistings % Add Quellcodeverzeichnis
\addcontentsline{toc}{section}{Quellcodeverzeichnis}

% das Abkürzungsverzeichnis
\newpage
% das Abkürzungsverzeichnis ausgeben
\fancyhead[L]{Abkürzungsverzeichnis} %Kopfzeile links
\section*{Abkürzungsverzeichnis}

\begin{acronym}
\end{acronym}

% \printnomenclature[3cm]
% Abkürzungsverzeichnis soll im Inhaltsverzeichnis auftauchen
\addcontentsline{toc}{section}{Abkürzungsverzeichnis}


%%%%%%% EINLEITUNG %%%%%%%%%%%%
\newpage
\fancyhead[L]{\nouppercase{\leftmark}} %Kopfzeile links

% 1,5 facher Zeilenabstand
\onehalfspacing

% arabische Seitennummerierung ab hier
\pagenumbering{arabic}

% Alternative Einbindung des Abstract in Kapitel "0" falls gewünscht
%\setcounter{section}{-1}
%\setcounter{page}{0}

% Option: Einbindung abstract
%\section*{Kurzfassung}

Diese Arbeit analysiert die Echtzeitfähigkeit von Micro-ROS und FreeRTOS anhand
einer Robotersteuerungssoftware. Ziel ist der Vergleich beider Systeme
hinsichtlich ihres Echtzeitverhaltens mit Schwerpunkt auf den Ausführungszeiten.
Dazu wird zunächst die bestehende Micro-ROS-Steuerungssoftware auf FreeRTOS
portiert. Anschließend wird eine zyklengenaue Messmethode für den Programmlauf
basierend auf der Data Watchpoint and Trace Unit (DWT) der ARM-Architektur
entwickelt. Zum Abschluss werden die Laufzeitdaten visualisiert und ausgewertet.
Die Ergebnisse verdeutlichen die unterschiedlichen Schwerpunkte beider
Plattformen: Micro-ROS punktet durch die Kompatibilität mit dem ROS-Ökosystem,
während FreeRTOS mit minimalen Latenzen und deterministischem Scheduling
deutlich besseres Echtzeitverhalten bietet. Die Analyse bestätigt zudem den
maßgeblichen Einfluss der Cache-Nutzung auf die Rechenleistung.

\section*{Abstract}

This thesis analyzes the real-time capabilities of micro-ROS and FreeRTOS using
a robotic control software. The goal is to compare both systems regarding their
real-time behavior emphasized on the execution times. First, the existing
micro-ROS control software is ported to FreeRTOS. Next, a cycle-accurate
measurement method for program execution based on ARM's Data Watchpoint and
Trace Unit (DWT)is developed. Finally, the runtime data is evaluated. The
evaluation of the results highlights the different strengths of both platforms:
micro-ROS stands out due to its compatibility with the ROS ecosystem, while
FreeRTOS achieves superior real-time behavior through minimal latencies and
deterministic scheduling. Additionally, the analysis confirms the significant
impact of the cache utilization on system performance.

%\newpage

% einzelne Kapitel werden hier eingebunden
\section*{Einleitung}

Die vorliegende Arbeit beschäftigt sich mit der Analyse der Echtzeitfähigkeit
von Micro-ROS und FreeRTOS am Beispiel einer Robotersteuerungssoftware. Ziel ist
es, die Performance beider Systeme hinsichtlich der Ausführungszeiten,
Ressourcenverbrauch sowie Echtzeitverhalten zu untersuchen, um ihre Eignung für
Roboteranwendungen zu bewerten.

Die Arbeit beinhaltet schwerpunktmäßig die Entwicklung einer Methode zum
Profiling der Steuerungssoftware eines mobilen Roboters. Dabei wird zunächst die
bestehende Firmware, die auf Micro-ROS basiert, im Rahmen dieser Arbeit auf
FreeRTOS portiert. Anschließend werden die Methodik zur Generierung von
Profiling-Daten für die Analyse detailliert beschrieben und die resultierenden
Ergebnisse evaluiert.

Zu Beginn wird ein Überblick über die grundlegenden Konzepte gegeben.
Darauffolgend werden die Implementierungen detailliert vorgestellt. Abschließend
werden die erzielten Ergebnisse evaluiert, und es wird ein Ausblick auf weitere
Anwendungsmöglichkeiten und Optimierungspotenziale gegeben.

\newpage

\section{Hintergrund}

Die vorliegende Bachelorarbeit hat zum Ziel, die Robotersteuerungssoftware, die
derzeit auf Micro-ROS basiert, auf FreeRTOS zu portieren, um einen
vergleichenden Leistungsanalyse zwischen beiden Plattformen durchzuführen. Beide
Systeme sind für die Steuerung eines mobilen Roboters auf einem Cortex-M7
Mikrocontroller von Arm konzipiert, unterscheiden sich jedoch in ihrer
grundlegenden Architektur, was sich auch in ihrer Echtzeitfähigkeit und
Ressourcennutzung widerspiegelt. Während Micro-ROS auf der \ac{ROS}2 aufbaut und
eine höhere Abstraktionsebene sowie standardisierte Kommunikationsschnittstellen
mittels der \ac{DDS}-Middleware bietet, basiert Micro-ROS selbst auf FreeRTOS.
Die Portierung der Robotersteuerungssoftware von Micro-ROS auf FreeRTOS kann
daher als eine Reduzierung der Abhängigkeitsebene betrachtet werden. Dies
ermöglicht eine direktere und effizientere Nutzung der zugrunde liegenden
Echtzeit-, sowie Speicherressourcen.

\begin{figure}[htb]
    \centering
    \includegraphics[width=0.8\textwidth]{assets/Micro-ROS_architecture}
    \caption{Micro-ROS Architektur\cite[S. 6]{koubaa2023}}
% \label{fig:uros_architecture}
\end{figure}

Nach dem Wechsel zu FreeRTOS wird die Echtzeitleistung der Steuerungssoftware
analysiert mit einem besonderen Fokus auf den Overhead, der durch die
Micro-ROS-Schicht verursacht wird. Der Vergleich soll aufzeigen, inwiefern
FreeRTOS durch die Eliminierung dieser zusätzlichen Abhängigkeit eine
effizientere und leichtgewichtige Lösung für kritische Roboteranwendungen
darstellt. Dabei soll der Einsatz einer zyklengenaue Messung des Programmablaufs
ermöglichen, fundierte Aussagen über die Echtzeitfähigkeit beider Plattformen zu
treffen, und den Leistungsgewinn anhand von diesem Beispiel für eine
Steuerungssoftware quantitativ zu belegen.

\subsection{FreeRTOS}

% TODO entwickelt von wem

FreeRTOS ist ein leichtgewichtiges, Open-Source-Echtzeitbetriebssystem, das
speziell für eingebettete Systeme entwickelt wurde. Es zeichnet sich unter
anderem durch deterministisches Verhalten mit Echtzeitgarantie sowie
konfigurierbare Allokation Speicher aus. Diese Eigenschaften machen es zu einer
idealen Wahl für Robotersteuerungssoftware, insbesondere wenn
Echtzeitanforderungen und effiziente Ressourcennutzung im Vordergrund stehen.

Im Rahmen dieser Bachelorarbeit wird FreeRTOS als Alternative zu Micro-ROS
eingesetzt, um die bestehende Robotersteuerungssoftware zu portieren. Während
Micro-ROS auf der Robot Operating System (ROS) 2-Middleware basiert und eine
hohe Abstraktionsebene sowie standardisierte Kommunikationsschnittstellen
bietet, ist FreeRTOS darauf ausgelegt, deterministische Echtzeitanforderungen zu
erfüllen und kritische Aufgaben in eingebetteten Systemen effizient zu
verwalten.

Ziel der Migration ist es, die Echtzeitleistung beider Systeme zu vergleichen,
insbesondere in Bezug auf Latenzzeiten, Scheduling-Effizienz und
Speicherauslastung. FreeRTOS bietet hierfür geeignete Mechanismen wie Trace
Hooks, die eine zyklengenaue Messung der Systemleistung ermöglichen. Diese
Analyse soll aufzeigen, inwiefern FreeRTOS als Alternative zu Micro-ROS für
kritische Roboteranwendungen geeignet ist und welche Vorteile es in Bezug auf
Determinismus und Ressourceneffizienz bietet. \subsubsection{Konzepte und
Komponenten}

\subsubsection{Konzepte}

Hier ist die Sektion Konzepte und Komponenten unter Berücksichtigung der
gegebenen Anforderungen:

Konzepte und Komponenten FreeRTOS bietet eine Reihe von grundlegenden Konzepten
und Komponenten, die für die Entwicklung von eingebetteten Echtzeitsystemen
essenziell sind. Im Fokus dieser Arbeit stehen insbesondere Queues,
Semaphore/Mutexe, Direct Task Notifications und Trace Hooks, die für die
effiziente Steuerung und Synchronisation von Robotersteuerungssoftware von
zentraler Bedeutung sind.

\paragraph{Queues}

Queues sind eine der Kernkomponenten von FreeRTOS und dienen der
Interprozesskommunikation zwischen Tasks. Sie ermöglichen den sicheren Austausch
von Daten und Nachrichten, selbst in multithreaded Umgebungen. Queues sind
threadsicher und können sowohl zur Datenübertragung als auch zur
Task-Synchronisation verwendet werden. Viele der höheren
Synchronisationsmechanismen in FreeRTOS, wie Semaphore und Mutexe, sind auf
Queues aufgebaut, was ihre zentrale Rolle im System unterstreicht.


\paragraph{Semaphore und Mutexe}

Semaphore und Mutexe sind Synchronisationsmechanismen, die den Zugriff auf
gemeinsame Ressourcen koordinieren und Race Conditions verhindern.

Semaphore werden verwendet, um den Zugriff auf Ressourcen zu beschränken oder um
Ereignisse zwischen Tasks zu signalisieren. Mutexe (Mutual Exclusion) sind
spezialisierte Semaphore, die sicherstellen, dass nur ein Task gleichzeitig auf
eine kritische Ressource zugreifen kann.

\paragraph{Direct Task Notifications}

Direct Task Notifications sind ein effizienter und ressourcensparender
Mechanismus zur Task-Synchronisation. Im Gegensatz zu Semaphore oder Queues
senden sie direkte Signale an einen Task, ohne zusätzliche Datenstrukturen zu
benötigen. Dies reduziert den Overhead und ermöglicht eine schnellere
Synchronisation.

Vergleich mit Semaphore: Während Semaphore universell einsetzbar sind, bieten
Direct Task Notifications eine leistungsstarke Alternative, wenn nur einfache
Synchronisation zwischen Tasks erforderlich ist.

\paragraph{Trace Hooks}

Trace Hooks sind spezielle Vorhakpunkte in der FreeRTOS-API, die es ermöglichen,
Ereignisse im System zu verfolgen und zu protokollieren. Sie sind besonders
nützlich für die Echtzeitanalyse, da sie detaillierte Informationen über
Task-Wechsel, Scheduling und Synchronisationsereignisse liefern. Diese
Funktionen sind entscheidend, um die Leistungscharakteristika der
Robotersteuerungssoftware zu analysieren und Optimierungspotenziale zu
identifizieren.


Zusammenfassend bilden diese Komponenten die Grundlage für die effiziente und
deterministische Steuerung von Robotersteuerungssoftware in FreeRTOS. Sie
ermöglichen eine zuverlässige Kommunikation, Synchronisation und Analyse, die
für die erfolgreiche Migration und den Vergleich mit Micro-ROS unerlässlich
sind.

\subsection{Echtzeitanalyse}

Um die Echtzeitanalyse der Steuerungssoftware durchzuführen, ist eine Methode
erforderlich, die jeden Ausführungsabschnitt des Codes flexibel, präzise und
threadsicher messen kann. Da die Software multithreaded ist, muss sichergestellt
werden, dass die Messungen trotz Interrupts und preemptivem Scheduling korrekt
und zyklengetreu durchgeführt werden können.

\subsection{Zyklenzähler}

Basierend auf den oben genannten Herausforderungen bietet die \ac{DWT} eine
geeignete Lösung \cite{ARM_KA001499}. Die DWT ist ein Hardwaremodul in
Prozessoren inklusive ARMv7-M, das speziell für die Analyse von
Programmabläufen entwickelt wurde \cite{ARMv7_ref_man_dwt_profiling}. Ein für
diese Arbeit zentraler Bestandteil der DWT ist der Zyklenzähler
\mintinline{c}|DWT_CYCCNT|, der bei jedem Takt inkrementiert wird, solange sich
der Prozessor nicht im Debug-Zustand befindet \cite{ARMv7_ref_man_dwt_cycle}.
Dadurch ermöglicht die DWT beispielsweise die Erfassung von Echtzeitaspekten mit
zyklengenauer Präzision under normaler Operation \cite{ARMv7_ref_man_dwt}.

\subsubsection{Beispiel: SEGGER SystemView}

Ein Beispiel hierfür ist SEGGER SystemView, ein Echtzeit-Analysewerkzeug, das
den DWT-Zyklenzähler einsetzt, um Live-Code-Profiling auf eingebetteten Systemen
durchzuführen \cite{SEGGER_SystemView}.

% TODO maybe not inline code to bloat the thesis
Das SEGGER SystemView nutzt den DWT, indem die Funktion \linebreak
\mintinline{c}|SEGGER_SYSVIEW_GET_TIMESTAMP()| einfach die hardkodierte
Registeradresse des Zyklenzählers zurückgibt, anstatt die interne Funktion
\mintinline{c}|SEGGER_SYSVIEW_X_GetTimestamp()| aufzurufen
\cite{SystemView_Config, Arm_DWT_Programmers_Model}. Dies erfolgt, sofern das
Macro \mintinline{c}|SEGGER_SYSVIEW_CORE_CM3| definiert ist
\cite{SystemView_Config_79}. Aus dieser Definition wird gezeigt, dass die
Nutzung des DWT für Cortex-M3, M4 sowie M7 Prozessoren aktiviert wird. Daraus
könnte sich eventuell auch schlussfolgern lassen, dass der DWT nur auf diesen
Prozessoren verfügbar ist.

Aufgrund dieser Eigenschaften der DWT wird die Implementierung für die
Echtzeitanalyse auf dieser Komponente aufbauen.

\newpage

\section{Vorbereitung}

\subsection{Umstellung auf FreeRTOS}

lorem

\newpage

\section{Abschluss}

Am Anfang wurde die Robotersteuerungssoftware von Micro-ROS auf FreeRTOS
umgestellt. Anschließend wurden eine Multi-Producer-Senke sowie ein Verfahren
entwickelt, das Echtzeitinformationen über die Steuerungssoftware ausgeben kann.
Abschließend wurde die Echtzeitfähigkeit basierend auf die Echtzeitinformationen
analysiert.

\subsection{Fazit}

Es lässt sich schlussfolgern, dass die Steuerungssoftware zwar durch Integration
von Micro-ROS funktionsreicher und folglich mit einer Vielzahl von
ROS-Komponenten kompatibel wird, dies allerdings mit erheblichem Overhead
erkauft wird. Bei begrenztem Speicher oder Rechenleistung bleibt FreeRTOS mit
seinem schlanken Kernel und den threadsicheren Queue-Implementierungen weiterhin
eine geeignete Wahl gegenüber komplexeren RTOS-Lösungen -- besonders wenn harte
Echtzeitfähigkeit im Vordergrund steht.

Darüber hinaus wurde gezeigt, dass sich die Softwareleistung durch L1-Caches
deutlich verbessern lässt, was für leistungskritische Software entscheidend ist.

\subsection{Ausblick}

Für zukünftige Arbeiten könnte die Multi-Producer-Senke so weiterentwickelt
werden, dass sie atomare Schreiboperationen auf 4-Byte-/32-Bit-Ebene
unterstützt. Dadurch könnten die Echtzeit-Daten nicht mehr im menschenlesbaren
Format, sondern maschinenlesbar und komprimiert als 4-Byte-Dateneinheit
ausgegeben werden.

Dies würde den Zwischenpuffer für die Echtzeitanalyse überflüssig machen, da die
durch Kontextwelchsel/Interrupt generierten Daten direkt atomar in die Senke
geschrieben werden könnten. Dazu müsste ein Parser auf dem Host-System
entwickelt werden, der idealerweise auch die Visualisierung sowie Analyse der
Daten parallel in Echtzeit ermöglicht.

\newpage

% einfacher Zeilenabstand
\singlespacing
% Literaturliste soll im Inhaltsverzeichnis auftauchen
\newpage
\phantomsection
\addcontentsline{toc}{section}{Literaturverzeichnis}
% Literaturverzeichnis anzeigen
\renewcommand\refname{Literaturverzeichnis}
\bibliography{Hauptdatei}

%% Index soll Stichwortverzeichnis heissen
% \newpage
% % Stichwortverzeichnis soll im Inhaltsverzeichnis auftauchen
% \addcontentsline{toc}{section}{Stichwortverzeichnis}
% \renewcommand{\indexname}{Stichwortverzeichnis}
% % Stichwortverzeichnis endgültig anzeigen
% \printindex

% \onehalfspacing
% % evtl. Anhang
% \newpage
% \phantomsection
% \addcontentsline{toc}{section}{Anhang}
% \fancyhead[L]{Anhang} %Kopfzeile links
% \input{anhang}

\newpage
% \includepdf[pages=1, fitpaper=true]{eigenstaendigkeitserklaehrung.pdf}

% \fancyhead[L]{Eigenständigkeitserklährung} %Kopfzeile links
% \input{eigenstaendigkeitserklaehrung} % deleted 

\end{document}
